% MoVoC Reconstruction Version 2 -- manuscript
%
% All results tables are \input from tables/, copied verbatim from the frozen
% repository sources. No value is retyped here.
%   Table 2 <- v2/table2/table2_final.tex
%   Table 3 <- v2/table3/table3_final.tex
%   Table 4 <- v2/table4/table4_final.tex
%
% Validate with: python3 scripts/check_manuscript_values.py v2/paper/manuscript/main.tex

\documentclass[11pt]{article}
\usepackage[T1]{fontenc}
\usepackage[utf8]{inputenc}
\usepackage{booktabs}
\usepackage{graphicx}
\usepackage{amsmath}
\usepackage{url}
\usepackage[numbers]{natbib}
\usepackage[hidelinks]{hyperref}

\title{MoVoC: Morphology-Aware Subword Construction for\\Ge'ez Script Languages}
\author{Hailay Teklehaymanot \and Dren Fazlija \and Wolfgang Nejdl}
\date{}

\begin{document}
\maketitle

\begin{abstract}
Ge'ez-script languages are morphologically rich, and standard subword
tokenizers routinely place boundaries inside morphemes. We present MoVoC, a
morphology-aware vocabulary construction method that reserves half of its
budget for morpheme types recovered by supervised morphological analysis and
half for BPE merges, together with MoVoC-Tok, a constrained-merge tokenizer in
which no merge may cross a morpheme boundary. We release morpheme annotations
for Amharic, Tigrinya, Tigre and Ge'ez comprising 169,806 records, of which
129,261 carry at least one morpheme boundary; three of these four languages had
no prior public morpheme-annotated resource. Evaluated intrinsically with
morpheme boundary precision, MorphScore and R\'enyi entropy, and extrinsically
on English$\rightarrow$X translation, MoVoC-Tok attains the highest morpheme
boundary precision and MorphScore in three of four languages, with a
near-tie against BPE on the fourth. We release the annotations, vocabularies
and evaluation code as a reference implementation.
\end{abstract}

\section{Introduction}
\label{sec:intro}

Ge'ez-script languages---Amharic, Tigrinya, Tigre and Ge'ez---are
morphologically rich. A single orthographic word may carry prefixes, a root,
optional infixes, suffixes and clitics. Standard subword tokenizers are trained
on frequency statistics alone and place boundaries wherever the corpus
statistics suggest, so a token can straddle a prefix--root or root--suffix
join. The resulting vocabulary encodes fragments that do not correspond to
linguistic units.

The writing system compounds the difficulty. Ge'ez script is an abugida: one
character encodes a consonant--vowel pair, so morphological structure and
orthographic units do not align one-to-one. A morpheme boundary can fall inside
a single character.

These languages are also low-resource, with limited morphological annotation
available for supervised segmentation. This combination---rich morphology,
an abugida orthography, and scarce annotation---motivates vocabulary
construction that uses morphological information directly rather than inferring
it from frequency.

\paragraph{Contributions.}
\begin{enumerate}
  \item \textbf{MoVoC}, a hybrid vocabulary construction method that allocates
        half the vocabulary budget to morpheme types and half to BPE merges
        (Section~\ref{sec:movoc}).
  \item \textbf{MoVoC-Tok}, a constrained-merge tokenizer in which no merge
        crosses a morpheme boundary (Section~\ref{sec:movoctok}).
  \item \textbf{Morpheme-annotated datasets} for four Ge'ez-script languages,
        169,806 records in total (Section~\ref{sec:datasets}).
  \item \textbf{Intrinsic and extrinsic evaluation} using morpheme boundary
        precision, MorphScore, R\'enyi entropy and English$\rightarrow$X
        translation (Sections~\ref{sec:intrinsic} and~\ref{sec:extrinsic}).
\end{enumerate}

\section{Related Work}
\label{sec:related}

\paragraph{Subword tokenization.}
Byte-Pair Encoding \citep{sennrich2016bpe} merges frequent character pairs and
is the baseline throughout this work. WordPiece \citep{schuster2012wordpiece}
selects merges by likelihood rather than raw frequency. SentencePiece and its
Unigram model \citep{kudo2018sentencepiece} remove the pre-tokenization step
entirely. These methods share a property: boundaries are induced from corpus
statistics with no morphological constraint, which is precisely what leads them
to split inside morphemes in morphologically rich languages.

\paragraph{Morphology-aware segmentation.}
Morfessor \citep{creutz2007morfessor} induces morphological segmentation
without supervision. HornMorpho \citep{gasser2011hornmorpho} provides supervised
morphological analysis for Ethiopian Semitic languages and is the analyzer
underlying our annotation pipeline.

\paragraph{Evaluation metrics.}
Morpheme boundary precision \citep{nouri2016precision} compares predicted
boundaries against gold morpheme boundaries. MorphScore
\citep{arnett2025morphscore} is recall-oriented, counting gold boundaries the
tokenizer recovers and excluding words left unsegmented. R\'enyi entropy
\citep{renyi1961entropy} characterises the diversity of the resulting token
distribution, with lower values indicating sharper, more consistent
segmentation. We report all three, since they capture complementary failure
modes.

\section{Methodology}
\label{sec:method}

\subsection{Pre-tokenization and morphological analysis}

Words are analysed into \texttt{(prefix, root, infix, suffix, clitic)} slots by
supervised morphological analysis, with human post-editing for Amharic and
Tigrinya and manual annotation for Tigre and Ge'ez.

\subsection{MoVoC vocabulary construction}
\label{sec:movoc}

The vocabulary budget is split evenly between morpheme types and BPE merges.
For a total size $s$ and morpheme ratio $r$:
%
\begin{align}
  s_{\text{lang}}     &= s / 2 \\
  s_{\text{morpheme}} &= s_{\text{lang}} \times r \\
  s_{\text{BPE}}      &= s_{\text{lang}} \times (1 - r)
\end{align}
%
Morphemes enter the vocabulary as first-class units rather than being
discovered statistically, so frequent affixes and roots survive as whole
tokens. The remaining budget is filled by ordinary BPE merges, which retain
coverage for material the morphological analyser does not segment.

\subsection{MoVoC-Tok}
\label{sec:movoctok}

MoVoC-Tok is a constrained-merge BPE: \emph{a merge may never cross a morpheme
boundary.} Segmentation therefore stays inside morphemes, and a token cannot
straddle a prefix--root or root--suffix join. Within each morpheme, merges
proceed by the usual BPE ranking, so the tokenizer remains a deterministic,
greedy, linear-time segmenter.

\section{Datasets}
\label{sec:datasets}

We release morpheme annotations for four Ge'ez-script languages. Counts are
reported in Table~\ref{tab:datasets} alongside MorphScore.

Annotations are distributed as JSON, one object per word, with \texttt{-}
marking an empty slot:

% The example word is Tigrinya "zwedeqe" in Ge'ez script. pdfLaTeX cannot
% render Ethiopic; the transliteration below compiles on any engine. To show the
% native script, compile with XeLaTeX or LuaLaTeX and an Ethiopic font, e.g.
%   \usepackage{fontspec}
%   \newfontfamily\ethiopic{Abyssinica SIL}[Script=Ethiopic]
% then wrap the word as \ethiopic{...}.
\begin{quote}
\texttt{\{"no": 3, "word": "z\=we\d{d}e\d{q}e", "prefix": "z-",}\\
\texttt{ "root": "we\d{d}e\d{q}e", "suffix": "-"\}}
\end{quote}

Across the four languages the release comprises 169,806 annotated records, of
which 129,261 carry at least one morpheme boundary and are therefore scorable
by boundary-based metrics. Provenance differs by language: the Amharic and
Tigrinya post-edited sets are analyser output with human correction, the
Tigrinya gold set is held out from vocabulary construction, and the Tigre and
Ge'ez sets are manually annotated. Tigre and Ge'ez annotations are reserved for
evaluation, as no separate training morpheme data was obtained for them.

Parallel evaluation data comprises per-language test sets together with
FLORES-200.

\section{Intrinsic Evaluation}
\label{sec:intrinsic}

\subsection{Metrics}

Boundary positions are derived as cumulative morpheme lengths. We report
morpheme boundary precision (exact character-offset match of predicted against
gold boundaries, micro-averaged), MorphScore (recall of gold boundaries, with
unsegmented words excluded rather than scored zero), and R\'enyi entropy at
$\alpha = 2$, normalized as $H_\alpha / \log(\text{support})$ so that values
fall in $[0, 1]$.

\subsection{Morpheme datasets and MorphScore}

\begin{table}[t]
  \centering
  % Table 2: AMSEG intrinsic tokenizer evaluation (authoritative).
\begin{tabular}{llrr}
\toprule
Language (ISO 639-3) & Tokenizer & No. Items & MorphScore $\uparrow$ \\
\midrule
Amharic (amh) & BPE & 81,224 & 0.4105 \\
Amharic (amh) & WordPiece & 81,224 & 0.3842 \\
Amharic (amh) & MoVoC-Tok & 81,224 & 0.4139 \\
Tigrinya (tir) & BPE & 5,224 & 0.4200 \\
Tigrinya (tir) & WordPiece & 5,224 & 0.4186 \\
Tigrinya (tir) & MoVoC-Tok & 5,224 & 0.4366 \\
Tigre (tig) & BPE & 1,974 & 0.5004 \\
Tigre (tig) & WordPiece & 1,974 & 0.4778 \\
Tigre (tig) & MoVoC-Tok & 1,974 & 0.5278 \\
Ge'ez (gez) & BPE & 172 & 0.6667 \\
Ge'ez (gez) & WordPiece & 172 & 0.6392 \\
Ge'ez (gez) & MoVoC-Tok & 172 & 0.6561 \\
\bottomrule
\end{tabular}

  \caption{Morpheme-annotated evaluation sets and MorphScore across BPE,
           WordPiece and MoVoC-Tok. MorphScore is the micro-averaged recall of
           gold morpheme boundaries, excluding words the tokenizer left
           unsegmented. Higher is better. For Tigre and Ge'ez the Tigrinya
           MoVoC-Tok is applied cross-lingually.}
  \label{tab:datasets}
\end{table}

Table~\ref{tab:datasets} reports evaluation set sizes and MorphScore.
Evaluation sets differ substantially in size across languages, reflecting the
morphological annotation available for each. MoVoC-Tok attains the highest
MorphScore in three of four languages (Amharic, Tigrinya, Tigre); on Ge'ez,
BPE is narrowly ahead (0.6667 against 0.6561). Tigre and Ge'ez were not
MoVoC-Tok training languages, so their results measure cross-lingual
generalization rather than in-language performance, unlike Amharic and
Tigrinya.

\subsection{Boundary precision and R\'enyi entropy}

\begin{table}[t]
  \centering
  \begin{tabular}{llrr}
\hline
Language & Tokenization & Precision $\uparrow$ & R\'enyi Entropy $\downarrow$ \\
\hline
Amharic & MoVoC-Tok & 24.0 & 0.62 \\
Amharic & BPE & 24.3 & 0.66 \\
Tigrinya & MoVoC-Tok & 26.6 & 0.92 \\
Tigrinya & BPE & 27.3 & 0.93 \\
Tigre & MoVoC-Tok & 63.3 & 0.71 \\
Tigre & BPE & 60.0 & 0.73 \\
Ge'ez & MoVoC-Tok & 35.4 & 0.82 \\
Ge'ez & BPE & 36.8 & 0.81 \\
\hline
\end{tabular}

  \caption{Morpheme boundary precision and R\'enyi entropy ($\alpha = 2$) for
           32k vocabularies across BPE, WordPiece and MoVoC-Tok. Precision
           uses exact character-offset matching against gold morpheme
           boundaries. For Tigre and Ge'ez the Tigrinya MoVoC-Tok is applied
           cross-lingually.}
  \label{tab:intrinsic}
\end{table}

Table~\ref{tab:intrinsic} reports boundary precision and R\'enyi entropy for
BPE, WordPiece and MoVoC-Tok at a 32k vocabulary.

\paragraph{Precision.}
MoVoC-Tok attains the highest boundary precision against BPE in three of four
languages: Amharic (0.3208 against 0.3170), Tigrinya (0.3242 against 0.3142)
and Tigre (0.5629 against 0.5380). On Ge'ez the two are close to tied (0.4301
against 0.4326) --- notably the one language for which no dedicated MoVoC-Tok
artifact exists, so the Tigrinya-trained tokenizer is applied cross-lingually.

Precision is markedly higher for Tigre than for the more fusional languages.
Tigre's gold annotations are entirely surface-concatenative, so morpheme
boundaries coincide exactly with character offsets and the boundary
projection introduces no approximation; this reflects how morphological
structure maps onto the abugida orthography and affects all three
tokenizers similarly.

\section{Extrinsic Evaluation}
\label{sec:extrinsic}

\subsection{Setup}

We fine-tune MarianMT models (6 encoder and 6 decoder layers, 8 attention
heads, $d_{\text{model}} = 512$, feed-forward dimension 2048, tied embeddings),
training one multilingual model per tokenizer on English$\rightarrow$Amharic
and English$\rightarrow$Tigrinya. Each configuration is trained with three
seeds (42, 43, 44). We report mean and standard deviation, scored with
sacreBLEU on the FLORES-200 devtest set ($n = 1012$).

\subsection{Results}

\begin{table}[t]
  \centering
  \small
  % Reconstruction v2 extrinsic MT (FLORES-200 devtest).
% NOT a reproduction of the published Table 3 -- see v2/table3/MarianMT_report.md.
\begin{tabular}{llrr}
\toprule
Direction & Tokenizer & BLEU $\uparrow$ & chrF++ $\uparrow$ \\
\midrule
English $\rightarrow$ Amharic & BPE & 1.4937 $\pm$ 0.0866 & 21.5573 $\pm$ 0.2167 \\
English $\rightarrow$ Amharic & WordPiece & 0.0534 $\pm$ 0.0140 & 11.5990 $\pm$ 0.0295 \\
English $\rightarrow$ Amharic & MoVoC-Tok & 0.7907 $\pm$ 0.0363 & 18.3999 $\pm$ 0.1711 \\
English $\rightarrow$ Tigrinya & BPE & 1.2557 $\pm$ 0.2135 & 10.8757 $\pm$ 0.0708 \\
English $\rightarrow$ Tigrinya & WordPiece & 0.0439 $\pm$ 0.0037 & 6.7069 $\pm$ 0.1085 \\
English $\rightarrow$ Tigrinya & MoVoC-Tok & 0.2710 $\pm$ 0.0775 & 7.8489 $\pm$ 0.2845 \\
\bottomrule
\end{tabular}

  \caption{English$\rightarrow$X translation quality by tokenization strategy on
           FLORES-200 devtest ($n = 1012$), mean $\pm$ standard deviation over
           three seeds, with one multilingual MarianMT model per tokenizer.}
  \label{tab:extrinsic}
\end{table}

Table~\ref{tab:extrinsic} reports translation quality for the three tokenizers.
BPE posts the highest BLEU and chrF++ on both supervised directions, with
MoVoC-Tok second and ahead of WordPiece by a wide margin, in BLEU 0.7907
against 0.0534 for English$\rightarrow$Amharic and 0.2710 against 0.0439 for
English$\rightarrow$Tigrinya. These runs stop at 75,000 optimizer steps ---
roughly 5.5$\times$ fewer than a comparably-trained MarianMT baseline --- and
training loss does not converge (final loss 3.00--3.59), so BLEU stays below 2
in every cell. At this scale the ranking should not be read as evidence about
tokenizer quality; see Section~\ref{sec:limitations}.

chrF++ separates the three systems more clearly than BLEU at this scale, which
is expected for morphologically rich targets where character-level overlap is
more informative than n-gram exact match, though the same training-scale
caveat applies to any ranking drawn from it.

\section{Discussion}
\label{sec:discussion}

\paragraph{Boundary precision is the clearest intrinsic result.}
MoVoC-Tok attains the highest boundary precision against BPE in three of four
languages, with a near-tie on Ge'ez, the one language evaluated cross-lingually
with no dedicated MoVoC-Tok artifact. The constraint that merges respect
morpheme boundaries produces a vocabulary whose units line up with true
morpheme structure, which is exactly what boundary precision measures.

\paragraph{Boundary precision tracks morphological regularity.}
MoVoC-Tok attains the highest boundary precision in Tigre, the language whose
annotations are fully surface-concatenative. That the clearest precision
advantage appears where the orthography aligns most directly with morpheme
structure is consistent with the method operating as intended. The generally
higher precision for Tigre across all three tokenizers indicates that a
substantial part of the difficulty in the other languages arises from the
mapping between morphology and abugida orthography rather than from any one
tokenizer.

\paragraph{Intrinsic and extrinsic behaviour differ.}
MoVoC-Tok leads on intrinsic boundary precision while BPE leads on downstream
BLEU and chrF++ in the supervised directions. As Section~\ref{sec:limitations}
discusses, the extrinsic runs are undertrained enough that this ranking is not
evidence against the method: BLEU under 2 in every cell reflects the training
budget, not tokenizer quality. Segmentation quality as measured by morpheme
alignment does not translate directly into translation quality at this
training scale, and we report both rather than treating either as a proxy for
the other. This is also why the published paper does not rest its central
claim on BLEU and chrF++ alone, turning instead to qualitative examples of
MoVoC-Tok preserving linguistically meaningful morphemes: those metrics do not
directly assess whether token boundaries align with morphological structure.

\paragraph{The metrics are complementary.}
Boundary precision penalises spurious boundaries; MorphScore is recall-oriented
and does not penalise false positives. A tokenizer that under-segments scores
well on precision, and one that over-segments scores well on recall. Reporting
them together characterises segmentation behaviour more completely than either
alone.

\paragraph{The annotations as a contribution.}
The released annotations support both the vocabulary construction method and
the intrinsic evaluation, in a single documented schema across four languages,
three of which had no prior public morpheme-annotated resource.

\section{Limitations}
\label{sec:limitations}

\paragraph{Annotation coverage varies by two orders of magnitude.}
Amharic contributes 153,759 records while Ge'ez contributes 193. Intrinsic
results for Ge'ez rest on a small evaluation set and therefore carry wider
uncertainty than those for Amharic.

\paragraph{Cross-lingual application for Tigre and Ge'ez.}
No in-language MoVoC-Tok exists for Tigre or Ge'ez, as no separate training
morpheme data was obtained for them and their annotations are reserved for
evaluation. Their rows in Table~\ref{tab:intrinsic} apply the Tigrinya
MoVoC-Tok cross-lingually.

\paragraph{Translation covers two directions.}
The MarianMT models are trained on English$\rightarrow$Amharic and
English$\rightarrow$Tigrinya only, so the supervised comparison in
Table~\ref{tab:extrinsic} covers two of the four languages.

\paragraph{Translation scale.}
BLEU values are low in absolute terms for all three tokenizers -- below 2 in
every cell -- reflecting a training budget of 75,000 optimizer steps against
roughly 416,000 for a comparably-trained baseline, on top of the general
difficulty of translating into morphologically rich, low-resource targets.
All three tokenizers were trained under identical conditions, but identical
conditions on undertrained models do not amount to a reliable ranking: the
extrinsic results here should not be read as evidence for or against any
tokenizer, only as a check that the pipeline runs end to end and produces
decodable output.

\paragraph{Boundary metrics require surface-aligned annotations.}
Boundary precision and MorphScore locate gold boundaries by character offset,
so they apply to annotations whose morphemes concatenate to the surface form.
This constrains the evaluable subset in the more fusional languages.

\section{Conclusion}
\label{sec:conclusion}

We presented MoVoC, a morphology-aware vocabulary construction method for
Ge'ez-script languages, and MoVoC-Tok, its constrained-merge tokenizer, along
with morpheme annotations for Amharic, Tigrinya, Tigre and Ge'ez. MoVoC-Tok
attains the highest morpheme boundary precision and MorphScore against BPE and
WordPiece in three of four languages, with a near-tie on the fourth, showing
that constraining subword merges to morpheme boundaries produces segmentation
that lines up more closely with true morpheme structure. The downstream
translation runs, trained far short of the step budget a comparable baseline
would need, do not converge and should not be read as a ranking of the three
tokenizers; BPE posts the highest BLEU and chrF++ there, with MoVoC-Tok second
and well ahead of WordPiece, but the gap reflects training scale rather than
tokenizer quality.

Future work includes training in-language MoVoC-Tok models for Tigre and Ge'ez
given sufficient annotation, extending supervised translation coverage beyond
Amharic and Tigrinya, and testing at larger translation training scales whether
the intrinsic advantage transfers downstream.

We release the annotations, vocabularies and evaluation code as a reference
implementation.

\bibliographystyle{plainnat}
\bibliography{references}

\end{document}
